\newpage
\pagenumbering{arabic}
\setcounter{page}{1}

\chapter{Вовед}
\section{Мотивација}
Меките, деформабилни тела се голем дел од секојдневниот живот: ткаенина, мускули, кожа, перници, душеци, гуми, кал. Сепак, повеќе децении во компјутерската графика во реално време (real-time CG) овие тела биле третирани како графички трикови наместо физички објекти. За претставување на ткаенина во компјутерски анимации се правеле трикови со vertex shaders и секоја деформација била мануелно креирарана од аниматор кој одел слика по слика за да добие прифатлив резултат. Преминувањето кон физичкото симулирање на меки тела е постепена, но сега е се користи широко.

Првите техники за симулации на меки тела во реално време во продукција на анимација се на Провот \cite{Provot1995DeformationCI}, имплицитното интегрирање на ткаенина на Бараф и Виткин \cite{BaraffDavidWitkin1998} во 1998 и position-based dynamics на Милер \cite{pbd}. Секој од овие трудови додаваат нови можности за интерактивни деформации. Моменталната генерација, изградена врз XPBD (add cite), конечно дозволува вакви симулации со интерактивна брзина. 

Во видео игрите е зголемена побарувачката за интерактивност со светот. Модерните играчки погони?? како Unreal Engine 5, Unity и Frostbite доаѓаат со имплементирани симулации на ткаенина; игрите почесто вклучуваат интеракции со меки тела (деформабилни терени, органски околини, оштетување на каросерија во тркачки игри). Алатки како NVIDIA PhysX, модулот за меки тела на Bullet и Havok Cloth ги направија меките тела честа но се уште скапа одлика.

Во виртуелната реалност, побарувачката оди подалеку. Кога корисникот ќе фати деформабилен објект со неговата рака, секој визуелен артефакт е повидлив од обично: тунелирање, нестабилност, тресење или нереалистично деформирање сите ја кршат илузијата. Виртуелната реалност ги обновува интересите за робусни техники за детекција на колизија токму заради тоа што грешките се толку видливи - некоја грешка што корисникот би ја толерирал на монитор станува проблем кога неговата рака е во објектот.? Во медицинско и хируршко тренирање, меките тела не се специјален ефект туку се главна грижа. Хируршките симулатори (Simbionix, Mimic Technologies, академски системи од SOFA) имаат потреба од реалистична симулација на меки ткива, крвни садови и кожа кога се под контакт на хируршки алати. Кај нив грешните симулации може да го нарушат тренингот и да направат погрешни навики кај ученикот. 

Во анимацијата, филмовите и специјалните ефекти, каде што најчесто симулациите се извршуваат офлајн (не се во реално време) се применуваат интерактивни симулации на меки тела во фазите на предпродукција. Houdini, Blender и SideFX имаат симулации кои ја замаглуваат границата помеѓу предпродукциски интерактивни симулации и финалниот продукт. Во интернет продавниците и виртуелни системи за испробување облека симулаторите треба да „облечат“ произволно 3Д човечко тело со ткаенини што точно се спуштаат на телото. Ова е прилично нова област во која се користат СМТ (симулации на меки тела).

\section{Проблемот на колизии за деформабилни тела}
За да се разбере зошто проблемот на колизија на меки тела е потежок од колизијата кај тврдите тела, корисно е да се почне од тврдите тела. Тврдиот објект има фиксна геометрија: кутија, топка, мрежа од триаголници итн. Проверките за растојание до геометријата може да се пресмета однапред (со Signed Distance Fields, BVH или GJK/EPA) и контактот помеѓу две тврди тела се сведува на манипулирање на конечни контактни многобразија - најчесето точки, рабови или страни. Геометријата не се изменува во векторскиот простор на телото, па било која структура за забрзување која што се гради еднаш останува валидна низ целиот век на симулацијата. Децениите истражувања имаат произведено мноштво оптимизирани алгоритми за решавање на овој проблем.

Меките тели од другастрана ги прекршуваат сите претпоставки врз кои се потпираат техниките за колизии на тврди тела. Во остатокот од овој дел ќе се набројат главните проблеми, а главниот дел од дипломската работа (5ти и 6ти дел) ќе се фокусира на техниките кои се користат за да се решат проблемите. 

\textbf{Геометријата се променува на секоја слика(frame)}. „Обликот“ на мекото тело е моменталната конфигурација на точките од кои е составено. Тој облик се менува на секој временски чекор на симулацијата. Било која структура за пресметка на растојанието (BVH, SDF) која што е пресметана за телото мора да се изгради одново на секој временски чекор. Цената на реизградбата станува една од главните грижи, дрвото што е изградено еднаш на почеток и не се менува веќе го нема.

\textbf{Бројот на контакти е огромен}. Репрезентацијата на меки тела која што е мрежа од стотици или илјадници точки може да има истовремен контакт со друг објект кај десетици или стотици од овие точки. Множеството на контактни точки не е само мал дел од објектот туку цело множество, па алгоритмите што лошо скалираат со бројот на контакти веднаш наоѓаат на проблеми. Ова изнудува добри избори за податочни структури и алгоритми што овозможуваат паралалелно процесирање и што се добри за кеширање.

\textbf{Колизии на телото само со себе се тежок проблем}. Мекото тело може да се превитка и да се судри само со себе. Овие сопствени колизии значат дека секоја точка потенцијално може да се судри со секоја друга точка од телото, што за пресметување е $O(n^2)$ без структури за забрзување. Широката фаза на детекција на колизии треба да биде ефикасна и тела што се движат. Униформно хеширање (Teschner et al 2003) е најчесто решение и се користи во овој труд.
	

\textbf{Тунелирање}. Кога брзината на една точка во еден временски чекор е поголема од дебелината на некој објект со кој треба да се судри, може таа точка да помине низ објектот без колизијата да се детектира. Ова не е проблем уникатен за меки тела, но е почеста појава бидејќи повеќе честички се движат одеднаш, и „лесни“ честички кои добиваат големо забрзување може моментално да имаат високи брзини. Континуирана детекција на колизии (CCD) го решава овој проблем но со голема цена на пресметка, во овој труд се користи само дискретна детекција и се спомнува ова ограничување.

\textbf{Геометриски дисконтинуитети кај одликите на мрежата на објектот}. Кога се прават колизии на честички со некој произволен објект, на??? Види покасно 

Сите овие проблеми се генерално \textit{ортогонални}. Решавањето на еден од нив не ги решава останатите. Точен систем за симулација на меки тела треба да има множество од алатки кои работат заедно наместо само еден алгоритам. Мапирањето на тие алатки е целта на овој труд.

\section{Цели и преглед на трудот}

Една цел на овој труд е да се проучат главните техники за симулирање на колизии во модерните симулатори на меки тела - забрзување на широк појас, забрзување на тесен појас, сопствени колизии и реакции. Втората цел е да се демонстрира имплементациија на овие техники во симулатор на меки тела во реално време базиран на XPBD и напишан во C++. 

Остатокот од документот е структуриран на следниот начин:

\textit{Глава 2 (Позадина и поврзани трудови)} дава општ преглед на методите на симулација на меки тела: системи на маса и еластична пружина, методи на конечни елементи, динамика базирана на позиции. Потоа ги проучува пристапите кои се користат во постоечките физички симулатори и трудовите на кои се базира овој дипломски труд.

\textit{Глава 3 (Математички основи)} ги етаблира математичките методи: Њутново движење и Ојлерова интеграција, динамика базирана на ограничувања (constraint-based dynamics) со Лагранжови множители, формулацијата на XPBD (extended position-based dynamics). Исто така се воведуваат геометриските примитиви што се појавуваат во останатите глави.

\textit{Глава 4 (Модел на меко тело)} ја објаснува дискретизацијата: податочната структура што се користи, методот за тетраедрализација на објектот, разликата помеѓу внатрешни и надворешни честички и формулацијата на ограничувањата за рабови и волумен. 

\textit{Глава 5 (Техники за детекција на колизии)} е главниот дел. Се преставуваат забрзување во широка фаза(BVH и просторно кеширање), забрзување на тесна фаза кај мрежи од триаголници(најблиска точка, псевдо-нормали) и сопствени колизии. Секоја подглава алтернира помеѓу генералната техника и конкретната имплементација.

\textit{Глава 6 (Техники за реакција на колизии)} 

\textit{Глава 7 (Имплементација и рендерирање)} ја покрива имплементацијата во код на симулаторот: структурата на проектот, rendering pipeline, контроли со корисничкиот интерфејс.

\textit{Глава 8 (Резултати и евалуација)} претставува тест сценарија и стрес тестирање на симулацијата. 

\textit{Глава 9 (Заклучок)} се сумаризира направеното и се дава осврт на идни подобрувања: непрекинато детектирање на колизии, колизии помеѓу 2 меки тела, паралелизација, додавање материјали.


